\documentclass[spanish,letterpaper,oneside]{article}

\def\documenttitle {Calibrar para Decidir: Proyecto Integrador de Modelación Matemática}
\def\documentsubtitle {Fundamentos para la enseñanza y el aprendizaje I · Cohorte SD26}
\def\documentsubject {Actividad 8 · Proyecto Integrador Final}
\def\documentauthor {Mtro. Martín Jonathan de la Cruz Muñoz}
\def\coursename {Fundamentos para la enseñanza y el aprendizaje I}
\def\coursecode {MAEM-SD26}

\def\universityname {Instituto de Investigación, Innovación y Estudios de Posgrado para la Educación del Estado de Nuevo León}
\def\universityfaculty {Maestría en Aprendizaje y Enseñanza de las Matemáticas}
\def\universitydepartment {Fundamentos para la enseñanza y el aprendizaje I}
\def\universitydepartmentimage {img/iiiepe/marca-iiiepe.png}
\def\universitydepartmentimagecfg {height=1.75cm}
\def\universitylocation {Monterrey, Nuevo León}
\def\studentfullname {Mtro. Martín Jonathan de la Cruz Muñoz}
\def\studentid {00874}
\def\studentemail {martin.delacruz@campus.iiiepe.edu.mx}
\def\programperiod {2do periodo, septiembre--diciembre de 2026 (SD26)}
\def\programmodality {Modalidad presencial}
\def\programshort {MAEM}

\def\authortable {
  \begin{tabular}{ll}
    \textbf{Alumno:} & \studentfullname \\
    Matrícula: & \studentid \\
    Correo institucional: & \studentemail \\
    Unidad: & \coursename \\
    Programa: & \programshort \\
    Periodo: & \programperiod \\
    & \\
    \multicolumn{2}{l}{Fecha de realización: 19 de noviembre de 2026} \\
    \multicolumn{2}{l}{\universitylocation}
  \end{tabular}
}

\input{template}

\usepackage{helvet}
\renewcommand{\familydefault}{\sfdefault}
\usepackage{array}
\usepackage{longtable}
\usepackage{pdflscape}
\usepackage{tikz}
\usepackage{pgfplots}
\pgfplotsset{compat=1.18}
\usepackage[most]{tcolorbox}
\usetikzlibrary{arrows.meta,positioning,calc,fit,backgrounds,shapes.geometric}

\definecolor{iiiepeNavy}{HTML}{163B5C}
\definecolor{iiiepeBlue}{HTML}{1F6F9F}
\definecolor{iiiepeAqua}{HTML}{20B7A8}
\definecolor{iiiepeOrange}{HTML}{F28C28}
\definecolor{iiiepePaper}{HTML}{F6F8F7}
\definecolor{iiiepeInk}{HTML}{1E2A32}
\definecolor{iiiepePale}{HTML}{EEF5F8}
\definecolor{iiiepeMint}{HTML}{E7F6F4}
\definecolor{iiiepeWarm}{HTML}{FFF3E6}
\def\portraittitlecolor {iiiepeNavy}
\setcitestyle{authoryear,open={(},close={)}}
\makeatletter
\renewcommand{\@biblabel}[1]{}
\makeatother
\setlength{\bibhang}{1.27cm}
\setlength{\bibsep}{6pt}
\AtBeginDocument{\renewcommand{\refname}{Referencias}}

\newcommand{\insertcoverwatermark}{%
  \AddToShipoutPictureBG*{%
    \begin{tikzpicture}[remember picture,overlay]
      \node[opacity=0.04,inner sep=0pt] at (current page.center) {%
        \includegraphics[width=0.65\paperwidth]{img/iiiepe/marca-iiiepe.png}%
      };
    \end{tikzpicture}%
  }%
}

\setlength{\parskip}{5pt}\setlength{\parindent}{0pt}
\begin{document}\insertcoverwatermark\templatePortrait\templatePagecfg\templateFinalcfg\fancyfoot{}\fancyfoot[L]{\scriptsize IIIEPE · MAEM · SD26}\fancyfoot[C]{\scriptsize\thepage}\fancyfoot[R]{\scriptsize Actividad 8 · Trabajo Final}\setstretch{1.3}\raggedbottom
\section{1. Portada y planteamiento integrador del proyecto}
El presente proyecto integrador desarrolla de manera articulada el diseño completo de un \textbf{proyecto educativo por realizar}: \textit{«Calibrar para decidir: modelación matemática, aprendizaje profundo y uso responsable de sistemas de medición en la formación técnica»}. Mi postura es que la enseñanza de las matemáticas en bachillerato técnico e ingeniería alcanza pleno sentido formativo cuando articula el campo \textit{Saberes y Pensamiento Científico} \citep{sep2022plan,sep2024fase6}, una metodología sociocrítica de \textit{Aprendizaje Basado en Indagación (STEAM)} y la comprensión contextual del \textit{Modelo Bioecológico} \citep{bronfenbrenner2006bioecologico}.

Considero indispensable señalar de forma explícita y honesta que este proyecto corresponde a una \textbf{propuesta didáctica por realizar} y que los datos numéricos de calibración presentados ($y = 1.05x + 0.2$) tienen carácter \textbf{estrictamente didáctico y simulado}; en ningún caso se reportan resultados empíricos previos ni mediciones atribuidas falsamente a un grupo escolar. El estudio desarrolla diez apartados estructurados de fundamentación, metodología y prospectiva, e incorpora íntegramente la matriz de autoevaluación de veinte indicadores y reflexión metacognitiva.

\section{2. Descripción general del proyecto}
\begin{itemize}
  \item \textbf{¿En qué consiste el proyecto?} En el diseño de una secuencia didáctica de indagación cuantitativa donde equipos de estudiantes analizan mediciones reales o simuladas de un instrumento de taller (termómetro digital, báscula o galga de banco) frente a lecturas patrón, ajustan el modelo matemático que describe el comportamiento del sensor, evalúan su error sistemático y dictaminan cuándo resulta confiable tomar una decisión técnica a partir de dichas lecturas.
  \item \textbf{¿A qué nivel educativo está dirigido?} A estudiantes de educación media superior tecnológica en las especialidades de mantenimiento industrial y mecatrónica (CONALEP y CECyTE, cuartos a sextos semestres), siendo directamente adaptable a primeros semestres universitarios de Ingeniería Mecánica, de Sistemas e Industrial \citep{conalep2025mecatronica}.
  \item \textbf{¿Cuál es su propósito?} Formar en el estudiantado la capacidad de construir, contrastar y justificar modelos lineales con sentido físico y estadístico, integrando interpretación crítica de datos, trabajo colaborativo equitativo y toma de decisiones responsables.
  \item \textbf{¿Qué contenidos matemáticos desarrolla?} Proporcionalidad directa frente a relaciones lineales afines ($y = kx$ frente a $y = mx + b$), tasa de cambio y sesgo inicial, representaciones tabular, gráfica y algebraica, error absoluto y relativo, medidas de tendencia central y dispersión básica, análisis dimensional y validez de dominio \citep{sep2023pensamiento,star2015algebra}.
\end{itemize}

\section{3. Problemática o necesidad educativa que atiende}
En talleres técnicos y asignaturas introductorias de ingeniería se observa con frecuencia que el alumnado ejecuta mecánicamente operaciones algebraicas —despejar incógnitas o aplicar la fórmula de la pendiente— sin cuestionar qué magnitud física miden los símbolos ni si un instrumento ofrece lecturas confiables antes de ajustar una máquina o reportar un diagnóstico. Esta brecha evidencia una enseñanza matemática fragmentada centrada en cálculo aislado \citep{sawyer2008optimizing}.

Abordar esta situación mediante el aprendizaje de las matemáticas resulta prioritario porque distinguir una variación proporcional limpia de una lectura con desviación inicial ($y = 1.05x + 0.2$), cuantificar márgenes de error y justificar un criterio de aceptación exige razonamiento algebraico, gráfico y estadístico en un escenario cotidiano del perfil técnico y ciudadano; sin embargo, la existencia real de esta dificultad en cada nuevo grupo deberá contrastarse primero mediante diagnóstico inicial antes de darla por supuesta.

\section{4. Campo formativo: Saberes y Pensamiento Científico (Matemáticas)}
En congruencia con el Plan de Estudio y los Programas Sintéticos \citep{sep2022plan,sep2024fase6}, el proyecto fortalece cinco ejes del pensamiento científico y matemático:
\begin{enumerate}
  \item \textbf{Resolución de problemas:} El grupo define variables relevantes (referencia $x$ y lectura $y$), formula conjeturas sobre la confiabilidad del instrumento y selecciona herramientas algebraicas o estadísticas pertinentes.
  \item \textbf{Pensamiento lógico y matemático:} Analiza invariantes y desviaciones (diferencias $y - x$ y cocientes $y/x$) para determinar si existe proporcionalidad o componente afín constante.
  \item \textbf{Razonamiento y argumentación:} Exige justificar por escrito los supuestos del modelo empleado, identificar fuentes posibles de error sistemático frente a ruido aleatorio y rebatir con evidencia afirmaciones infundadas \citep{eef2017mathematics}.
  \item \textbf{Modelación y representación matemática:} Conecta dinámicamente tablas de lecturas, gráficas cartesianas con banda de tolerancia y expresión simbólica (como ilustra la Figura~\ref{fig:act8-calibracion}).
  \item \textbf{Uso de las matemáticas en la vida cotidiana:} Traslada el análisis de confiabilidad instrumental a decisiones domésticas, de salud o de consumo responsable (básculas, medidores eléctricos, termómetros clínicos), vinculando rigor técnico y criterio ciudadano.
\end{enumerate}

\begin{figure}[htbp]
\centering
\begin{tikzpicture}
\begin{axis}[
  width=11.8cm, height=6.4cm,
  xlabel={Referencia patrón simulada $x$ (unidades de ensayo)},
  ylabel={Lectura observada $y$},
  xmin=0, xmax=10.5, ymin=0, ymax=11.5,
  grid=major, grid style={iiiepePale!95!gray},
  legend pos=north west,
  legend style={font=\scriptsize,draw=iiiepeNavy,fill=white},
  label style={font=\scriptsize}, tick label style={font=\scriptsize}
]
  \addplot[iiiepeNavy,dashed,thick,domain=0:10] {x};
  \addlegendentry{Referencia ideal: $y = x$}
  \addplot[iiiepeOrange!90!black,very thick,mark=*,mark options={fill=iiiepeOrange}] coordinates {
    (0,0.20) (2,2.30) (4,4.40) (6,6.50) (8,8.60) (10,10.70)
  };
  \addlegendentry{Ejemplo didáctico simulado: $y = 1.05x + 0.2$}
  \node[anchor=east,font=\scriptsize,fill=white,draw=iiiepeBlue,rounded corners=2pt,inner sep=3pt,align=center] at (axis cs:10.1,2.0)
    {Desviación: $y - x = 0.05x + 0.2$\\(crece linealmente con $x$)};
\end{axis}
\end{tikzpicture}
\caption{Ejemplo didáctico simulado de calibración: contraste entre proporcionalidad ideal ($y=x$) y modelo afín ($y=1.05x+0.2$).}
\label{fig:act8-calibracion}
\end{figure}

{\small\setstretch{1.15}\renewcommand{\arraystretch}{1.28}
\begin{longtable}{|c|c|c|c|p{0.44\linewidth}|}\hline
\rowcolor{iiiepePale} \textbf{Patrón $x$} & \textbf{Lectura $y$} & \textbf{Error $y-x$} & \textbf{Razón $y/x$} & \textbf{Interpretación matemática didáctica} \\\hline
\endfirsthead
\hline
\rowcolor{iiiepePale} \textbf{Patrón $x$} & \textbf{Lectura $y$} & \textbf{Error $y-x$} & \textbf{Razón $y/x$} & \textbf{Interpretación matemática didáctica} \\\hline
\endhead
0 & 0.20 & $+0.20$ & --- & Revela sesgo inicial (ordenada al origen $b=0.2 \neq 0$). \\\hline
2 & 2.30 & $+0.30$ & 1.150 & Descarta proporcionalidad pura al variar la razón $y/x$. \\\hline
4 & 4.40 & $+0.40$ & 1.100 & Incremento constante de $2.10$ por cada $\Delta x = 2 \Rightarrow m=1.05$. \\\hline
6 & 6.50 & $+0.50$ & 1.083 & El error absoluto crece linealmente con la magnitud medida. \\\hline
8 & 8.60 & $+0.60$ & 1.075 & Supera una tolerancia hipotética de $\pm 0.50$; requiere corrección. \\\hline
10 & 10.70 & $+0.70$ & 1.070 & Sustenta el dictamen técnico de recalibración o restricción de rango. \\\hline
\end{longtable}}

\section{5. Metodología sociocrítica de la Nueva Escuela Mexicana}
El proyecto adopta el \textbf{Aprendizaje Basado en Indagación bajo enfoque STEAM} como metodología sociocrítica de la NEM \citep{sep2022plan}. Esta elección se justifica porque la pregunta central —\textit{¿cuándo podemos confiar en una medición para tomar una decisión técnica responsable?}— demanda formular hipótesis previas, diseñar una prueba controlada, recolectar y organizar evidencia numérica, modelar matemáticamente las relaciones observadas y comunicar conclusiones verificables \citep{quiroga2022aprendizaje}. Asimismo, cada campo aporta una función genuina: la ciencia fundamenta la magnitud física y las fuentes de error; la tecnología organiza la captura digital o tabular sin reemplazar el juicio analítico; la ingeniería define tolerancias funcionales; las artes visuales optimizan la lectura crítica de gráficas; y las matemáticas estructuran la argumentación lógica \citep{rivas2017innovaciones}.

\section{6. Descripción de las principales actividades del proyecto}
{\small\setstretch{1.15}\renewcommand{\arraystretch}{1.3}
\begin{longtable}{|>{\raggedright\arraybackslash}p{0.18\linewidth}|>{\raggedright\arraybackslash}p{0.46\linewidth}|>{\raggedright\arraybackslash}p{0.30\linewidth}|}\hline
\rowcolor{iiiepePale} \textbf{Fase del proyecto} & \textbf{Acciones principales de los estudiantes} & \textbf{Evidencia formativa} \\\hline
\endfirsthead
\hline
\rowcolor{iiiepePale} \textbf{Fase del proyecto} & \textbf{Acciones principales de los estudiantes} & \textbf{Evidencia formativa} \\\hline
\endhead
1. Planteamiento del problema & Analizan un caso real de falla por lectura mal interpretada; registran su predicción individual inicial sobre la confiabilidad del sensor. & Nota diagnóstica de predicción previa. \\\hline
2. Investigación & Indagan el significado de calibración, sesgo de cero, error de escala y tolerancia permisible en contextos de taller o hogar. & Síntesis conceptual de términos metrológicos. \\\hline
3. Recolección y análisis de datos & Trabajan con lecturas reales de banco (o el conjunto simulado) organizando tablas de $x$, $y$, $y-x$ y cociente $y/x$. & Hoja de registro tabular verificada por pares. \\\hline
4. Resolución matemática & Trazan gráficas comparativas frente a $y=x$, calculan pendiente e intercepto e interpretan los límites de validez del modelo. & Desarrollo algebraico y gráfico anotado. \\\hline
5. Elaboración de productos & Redactan una ficha técnica de decisión indicando rango confiable, ecuación de corrección y advertencias de uso. & Dictamen técnico en una cuartilla con gráfica. \\\hline
6. Socialización de resultados & Exponen y defienden en plenaria su dictamen ante cuestionamientos de otros equipos sobre supuestos y márgenes de error. & Coevaluación oral argumentada. \\\hline
7. Evaluación formativa & Reciben retroalimentación descriptiva centrada en conceptos y entregan una versión revisada junto con reflexión metacognitiva \citep{eef2021feedback}. & Versión mejorada y matriz formativa. \\\hline
\end{longtable}}

\section{7. Relación con la perspectiva ecológica de Urie Bronfenbrenner}
De acuerdo con \citet{bronfenbrenner2006bioecologico}, \citet{bronfenbrenner1987ecologia} y \citet{tong2023bioecological}, el desarrollo de aprendizajes durante el proyecto depende de las condiciones ecológicas que envuelven los procesos proximales:
\begin{itemize}
  \item \textbf{Microsistema:} Interacciones directas entre docente, pares en equipo, equipo de taller y apoyo familiar; un clima respetuoso frente a dudas iniciales permite probar hipótesis sin ansiedad.
  \item \textbf{Mesosistema:} Articulación coherente entre la clase de matemáticas, los módulos técnicos de mantenimiento/mecatrónica y situaciones familiares de medición cotidiana.
  \item \textbf{Exosistema:} Disponibilidad institucional de instrumentos de medición, tiempo de clase frente a jornadas laborales de los estudiantes y acceso a calculadoras o software; el proyecto contempla hojas impresas y sets de datos replicables para evitar exclusión tecnológica.
  \item \textbf{Macrosistema:} Creencias sociales sobre quién «sirve para las matemáticas», cultura técnica orientada al cumplimiento rápido sin reflexión y políticas curriculares de la NEM \citep{unesco2015replantear}.
  \item \textbf{Cronosistema:} Evolución de las explicaciones desde la predicción diagnóstica inicial hasta el dictamen final revisado y su transferencia a unidades posteriores del semestre.
\end{itemize}

\section{8. Impacto esperado del proyecto}
Como proyecto por realizar, se espera promover cinco aprendizajes matemáticos observables tras su implementación: (a) distinguir con claridad relaciones proporcionales frente a funciones lineales afines en problemas auténticos; (b) interpretar pendiente, intercepto y error con sentido físico y dimensional; (c) fortalecer el trabajo colaborativo con roles alternos de argumentación; (d) consolidar pensamiento crítico frente a cifras aparentemente exactas pero sin respaldo contextual; y (e) comunicar decisiones técnicas claras sustentadas con evidencia gráfica y escrita \citep{bostrom2023formative}. Además, este enfoque formativo vincula la matemática con criterios ciudadanos; sin embargo, su alcance deberá contrastarse en aula mediante registros reales. Por ello, es decir, no obstante la coherencia teórica de la secuencia, ningún resultado de aprendizaje se da por verificado antes de recoger evidencias estudiantiles concretas.

\section{9. Reflexión final (cuatro preguntas orientadoras)}
\begin{enumerate}
  \item \textbf{¿Qué concepción del aprendizaje de las matemáticas se observa en el proyecto?} Se concibe aprender matemáticas como un proceso activo, relacional y acompañado de indagación y diálogo, nunca como simple repetición individual de despejes descontextualizados \citep{sawyer2008optimizing}.
  \item \textbf{¿Cómo contribuyen las metodologías sociocríticas de la NEM al desarrollo del pensamiento matemático?} La indagación STEAM conecta una necesidad técnica real con rigor modelador y participación comunitaria, obligando a transformar fórmulas en argumentos útiles para decidir responsablemente \citep{sep2022plan}.
  \item \textbf{¿Qué aportaciones de la perspectiva ecológica de Bronfenbrenner se identifican en el proyecto?} Reconoce que los procesos proximales solo prosperan si se atienden accesibilidad material (exosistema), diálogo entre asignaturas (mesosistema) y trato digno libre de etiquetas deterministas en el aula (microsistema y macrosistema), orientando decisiones concretas de equidad sin suponer causalidades sin evidencia futura \citep{bronfenbrenner2006bioecologico}.
  \item \textbf{¿Qué beneficios aporta este tipo de proyectos a la formación integral del estudiante?} Integra mi base industrial de rigor metrológico y claridad estructural con una mirada pedagógica humanista que valora a cada estudiante en su contexto real, fortaleciendo autonomía intelectual, ética técnica y ciudadanía crítica.
\end{enumerate}

\section*{Autoevaluación del proyecto y reflexión personal}
Con el propósito de valorar reflexivamente el alcance formativo del proyecto integrador, se presenta la verificación sistemática de veinte indicadores de coherencia didáctica y dos preguntas de análisis metacognitivo:

{\scriptsize\setstretch{1.12}\renewcommand{\arraystretch}{1.22}
\begin{longtable}{|c|>{\raggedright\arraybackslash}p{0.62\linewidth}|c|c|c|}\hline
\rowcolor{iiiepePale} \textbf{\#} & \textbf{Indicador de autoevaluación del proyecto} & \textbf{Sí} & \textbf{En proceso} & \textbf{No} \\\hline
\endfirsthead
\hline
\rowcolor{iiiepePale} \textbf{\#} & \textbf{Indicador de autoevaluación del proyecto} & \textbf{Sí} & \textbf{En proceso} & \textbf{No} \\\hline
\endhead
1 & Elaboré una portada con el título del proyecto, mi nombre, la asignatura y la fecha. & X & & \\\hline
2 & Expliqué con claridad en qué consiste mi proyecto. & X & & \\\hline
3 & Mencioné el nivel educativo al que está dirigido. & X & & \\\hline
4 & Describí el propósito del proyecto. & X & & \\\hline
5 & Identifiqué los contenidos matemáticos que desarrolla. & X & & \\\hline
6 & Describí la problemática o necesidad que atiende mi proyecto. & X & & \\\hline
7 & Justifiqué por qué esta problemática es relevante. & X & & \\\hline
8 & Expliqué por qué es importante abordarla mediante el aprendizaje de las matemáticas. & X & & \\\hline
9 & Expliqué cómo mi proyecto desarrolla la resolución de problemas. & X & & \\\hline
10 & Expliqué cómo favorece el pensamiento lógico y matemático. & X & & \\\hline
11 & Expliqué cómo promueve el razonamiento y la argumentación. & X & & \\\hline
12 & Expliqué cómo favorece la modelación y representación matemática. & X & & \\\hline
13 & Relacioné las matemáticas con situaciones de la vida cotidiana. & X & & \\\hline
14 & Identifiqué la metodología utilizada y justifiqué mi elección. & X & & \\\hline
15 & Expliqué cómo esta metodología favorece el aprendizaje de las matemáticas. & X & & \\\hline
16 & Describí las actividades principales del proyecto. & X & & \\\hline
17 & Expliqué la relación del proyecto con los sistemas de Bronfenbrenner. & X & & \\\hline
18 & Expliqué los aprendizajes matemáticos esperados y su impacto. & X & & \\\hline
19 & Respondí de manera crítica las preguntas de reflexión final. & X & & \\\hline
20 & Incluí al menos tres referencias en formato APA 7. & X & & \\\hline
\end{longtable}}

\begin{tcolorbox}[colback=iiiepeMint,colframe=iiiepeAqua!80!black,title=\textbf{Reflexión personal sobre el diseño del proyecto},breakable]
\textbf{¿Qué fue lo que mejor logré en este proyecto?} Logré articular orgánicamente mi trayectoria técnica de ingeniería con los marcos pedagógicos trabajados a lo largo de la cohorte SD26 (ciencias del aprendizaje, aprendizaje profundo, investigación-acción y modelo bioecológico) dentro de un proyecto realista de modelación matemática de calibración, diferenciando con honestidad un diseño por implementar y datos didácticos simulados frente a evidencia empírica futura.

\medskip
\textbf{¿Qué aspecto necesito mejorar al implementar este proyecto?} Verificar al implementarse en grupos reales la calibración de tiempos destinados al trabajo colaborativo en equipos heterogéneos y recoger sistemáticamente los resultados del diagnóstico inicial para ajustar la complejidad de las guías escritas según los contextos particulares de cada grupo \citep{restrepo2004investigacion}.
\end{tcolorbox}

\clearpage
\section*{10. Transferencia profesional y conclusiones}
En el contexto laboral de bachillerato técnico e ingeniería, el proyecto integrado aquí diseñado consolida la transferencia coherente de toda la unidad de aprendizaje hacia una docencia humanista, contextualizada y disciplinariamente rigurosa. En conclusión, responder a los desafíos de la enseñanza matemática exige diseñar tareas con sentido real, propiciar procesos proximales equitativos y someter cada intervención a verificación reflexiva continua. Por tanto, los diez apartados desarrollados y la autoevaluación reflexiva consolidan una propuesta viable, éticamente fundada y lista para su implementación en aula.

\clearpage
\begin{thebibliography}{99}
\bibitem[Boström y Palm(2023)]{bostrom2023formative} Boström, E., y Palm, T. (2023). Formative assessment practice in mathematics: Teachers' perceptions of prerequisite conditions. \textit{Frontiers in Education}, 8, Artículo 1101192.
\bibitem[Bronfenbrenner(1987)]{bronfenbrenner1987ecologia} Bronfenbrenner, U. (1987). \textit{La ecología del desarrollo humano}. Paidós.
\bibitem[Bronfenbrenner y Morris(2006)]{bronfenbrenner2006bioecologico} Bronfenbrenner, U., y Morris, P. A. (2006). The bioecological model of human development. En R. M. Lerner (Ed.), \textit{Handbook of child psychology: Vol. 1. Theoretical models of human development} (6.ª ed., pp. 793--828). John Wiley \& Sons.
\bibitem[Collin y Quigley(2021)]{eef2021feedback} Collin, J., y Quigley, A. (2021). \textit{Teacher feedback to improve pupil learning}. Education Endowment Foundation.
\bibitem[CONALEP(2025)]{conalep2025mecatronica} Colegio Nacional de Educación Profesional Técnica. (2025). \textit{Perfil de egreso: Profesional Técnico-Bachiller en Mecatrónica. Modelo Académico 2025}. Secretaría Académica.
\bibitem[Henderson et~al.(2017)]{eef2017mathematics} Henderson, P., Hodgen, J., Foster, C., y Kuchemann, D. (2017). \textit{Improving mathematics in Key Stages 2 and 3}. Education Endowment Foundation.
\bibitem[Quiroga Aguilar y Lara Órdenes(2022)]{quiroga2022aprendizaje} Quiroga Aguilar, L., y Lara Órdenes, E. (2022). El aprendizaje profundo como herramienta para cambio en la visión de aprendizaje de una cultura escolar. \textit{Revista Educación las Américas}, 12(1).
\bibitem[Restrepo Gómez(2004)]{restrepo2004investigacion} Restrepo Gómez, B. (2004). La investigación-acción educativa y la construcción de saber pedagógico. \textit{Educación y Educadores}, (7), 45--55.
\bibitem[Rivas et~al.(2017)]{rivas2017innovaciones} Rivas, A., André, F., y Delgado, L. E. (Eds.). (2017). \textit{50 innovaciones educativas para escuelas}. CIPPEC / Santillana.
\bibitem[Sawyer(2008)]{sawyer2008optimizing} Sawyer, R. K. (2008). Optimizing learning: Implications of learning sciences research. OECD Publishing.
\bibitem[SEP(2022)]{sep2022plan} Secretaría de Educación Pública. (2022). \textit{Plan de Estudio para la educación preescolar, primaria y secundaria}. SEP.
\bibitem[SEP(2023)]{sep2023pensamiento} Secretaría de Educación Pública. (2023). \textit{Programa de estudios del recurso sociocognitivo Pensamiento Matemático I}. Subsecretaría de Educación Media Superior.
\bibitem[SEP(2024)]{sep2024fase6} Secretaría de Educación Pública. (2024). \textit{Programa de estudio para la educación secundaria: Programa sintético de la Fase 6}. SEP.
\bibitem[Star et~al.(2015)]{star2015algebra} Star, J. R., Caronongan, P., Foegen, A., Furgeson, J., Keating, B., Larson, M. R., Lyskawa, J., McCallum, W. G., Porath, J., y Zbiek, R. M. (2015). \textit{Teaching strategies for improving algebra knowledge}. IES.
\bibitem[Tong y An(2023)]{tong2023bioecological} Tong, X., y An, X. (2023). A review of empirical studies applying Bronfenbrenner's bioecological theory. \textit{Frontiers in Psychology}, 14, Artículo 1233925.
\bibitem[UNESCO(2015)]{unesco2015replantear} UNESCO. (2015). \textit{Replantear la educación: ¿Hacia un bien común mundial?} Ediciones UNESCO.
\end{thebibliography}

\end{document}
